% AA-32 | Iteracja XVI; szczegolowy audyt historyczny IV
\chapter{Wyróżnik wielomianu i grupa Galois}\label{aa:ch:discriminant-galois}
\section{Wyróżnik i znak permutacji}

Jeżeli wielomian ma dwa równe pierwiastki, różnica między
nimi jest zerem. Chcemy otrzymać z tych różnic jedno
wyrażenie, które nie zależy od numeracji pierwiastków
i daje się obliczać ze współczynników wielomianu.
Zastosowanie do grup Galois wymaga starannego rozróżnienia
charakterystyk: w charakterystyce różnej od $2$
wyróżnik wykrywa, czy automorfizmy Galois działają
jako permutacje parzyste; w charakterystyce $2$
wniosek tego rodzaju na ogół jest fałszywy.

\begin{definicja}\label{aa:def:polynomial-discriminant}
Niech $K$ będzie ciałem, $n\ge1$ liczbą całkowitą,
$f\in K[X]$ monicznym wielomianem stopnia $n$,
a $\overline K$ ustalonym domknięciem algebraicznym $K$.
Niech $\alpha_1,\ldots,\alpha_n\in\overline K$
będą wszystkimi pierwiastkami $f$, wypisanymi
z krotnościami, czyli
$f(X)=\prod_{j=1}^n(X-\alpha_j)$ w $\overline K[X]$.
\emph{Wyróżnikiem} $f$ nazywamy element
\[
 \operatorname{disc}(f)=
 \prod_{1\le i<j\le n}(\alpha_i-\alpha_j)^2\in\overline K.
\]
Przy $n=1$ iloczyn pusty ma wartość $1$.
\end{definicja}

\begin{przyklad}[Stopień drugi]
Dla $f(X)=X^2+bX+c$ i pierwiastków $\alpha,\beta$
zachodzą $\alpha+\beta=-b$ oraz $\alpha\beta=c$.
Stąd
\[
 \operatorname{disc}(f)=(\alpha-\beta)^2
 =(\alpha+\beta)^2-4\alpha\beta=b^2-4c.
\]
Nie potrzebujemy wzoru na pierwiastki; rachunek jest
poprawny nad dowolnym ciałem, również w charakterystyce $2$.
\end{przyklad}

\begin{twierdzenie}[Wyróżnik a rozdzielczość]\label{aa:tw:disc-separable}
Dla dowolnego ciała $K$ i każdego monicznego
$f\in K[X]$ stopnia dodatniego zachodzą dwa stwierdzenia:
$\operatorname{disc}(f)\in K$, a ponadto
$\operatorname{disc}(f)\ne0$ wtedy i tylko wtedy,
gdy $f$ jest rozdzielczy nad $K$, to znaczy ma
$\deg f$ różnych pierwiastków w $\overline K$.
\end{twierdzenie}
\begin{proof}
\textbf{Przynależność do $K$.}
Wyrażenie $D(X_1,\ldots,X_n)=
\prod_{i<j}(X_i-X_j)^2$ jest symetryczne:
permutacja zmiennych permutuje czynniki z możliwą
zmianą znaku różnic, lecz znak znika po podniesieniu
do kwadratu. Ponadto $D$ ma współczynniki całkowite,
więc po przejściu do $K$ należy do
$K[X_1,\ldots,X_n]^{S_n}$.
Twierdzenie~\ref{aa:tw:fundamental-symmetric} daje
$D=P(e_1,\ldots,e_n)$ dla pewnego $P\in K[T_1,\ldots,T_n]$.
Podstawiając pierwiastki i korzystając z
Twierdzenia~\ref{aa:tw:vieta-general}, dostajemy
\[
 \operatorname{disc}(f)
 =P(-c_1,c_2,\ldots,(-1)^nc_n)\in K,
\]
gdzie $f(X)=X^n+c_1X^{n-1}+\cdots+c_n$.

\textbf{Kryterium niezerowości.}
W $\overline K$, które jest ciałem, iloczyn
niezerowych elementów jest niezerowy. Wobec tego
$\operatorname{disc}(f)\ne0$ dokładnie wtedy,
gdy dla wszystkich $1\le i<j\le n$ mamy
$\alpha_i-\alpha_j\ne0$. Oznacza to, że wszystkie
wypisane pierwiastki są parami różne, czyli że
$f$ jest rozdzielczy. Argument działa dla każdej
charakterystyki i obejmuje wielomiany stopnia $1$.
\end{proof}

\begin{przyklad}[Wyróżnik zerowy w charakterystyce dwa]\label{aa:prz:disc-inseparable-char2}
Niech $K=\mathbb F_2(t)$, gdzie $t$ jest elementem
przestępnym nad $\mathbb F_2$, i niech $f(X)=X^2+t$.
W domknięciu algebraicznym $K$ istnieje element $\alpha$
taki, że $\alpha^2=t$. Ponieważ charakterystyka wynosi $2$,
mamy $(X-\alpha)^2=X^2+\alpha^2=X^2+t$.
Wielomian ma zatem jeden pierwiastek z krotnością $2$,
a jego wyróżnik wynosi $0$.
Z drugiej strony $f$ jest nierozkładalny nad $K$:
gdyby $t=(A(t)/B(t))^2$ dla względnie pierwszych
niezerowych wielomianów $A,B\in\mathbb F_2[t]$, to
$2\deg A=1+2\deg B$, co jest niemożliwe.
Przykład rozróżnia nierozkładalność od rozdzielczości:
w dodatniej charakterystyce nierozkładalny wielomian
może mieć wielokrotny pierwiastek.
\end{przyklad}

\begin{twierdzenie}[Kryterium grupy alternującej]\label{aa:tw:disc-alternating}
Niech $K$ będzie ciałem charakterystyki różnej od $2$,
$f\in K[X]$ monicznym i rozdzielczym wielomianem
stopnia $n\ge2$, $L$ jego ciałem rozkładu oraz
$G=\operatorname{Gal}(L/K)$.
Ustalmy numerację różnych pierwiastków; daje ona wierne
zanurzenie $G\hookrightarrow S_n$, ponieważ element $G$
jest określony przez obrazy pierwiastków generujących $L$.
Wówczas $G\subseteq A_n$ wtedy i tylko wtedy,
gdy $\operatorname{disc}(f)$ jest kwadratem w $K$.
\end{twierdzenie}
\begin{proof}
Wybieramy kolejność $\alpha_1,\ldots,\alpha_n$ i kładziemy
\[
 \Delta=\prod_{1\le i<j\le n}(\alpha_i-\alpha_j)\in L.
\]
Z rozdzielczości wynika $\Delta\ne0$, a definicja
daje $\Delta^2=\operatorname{disc}(f)$.
Dla dowolnej transpozycji dwóch indeksów
w iloczynie $\Delta$ znak zmienia się dokładnie
raz: czynnik różnicy zamienianych pierwiastków
zmienia znak, natomiast pozostałe czynniki
zostają odpowiednio poprzestawiane parami.
Każdą permutację można przedstawić jako iloczyn
transpozycji, więc dla każdego $\sigma\in G$ zachodzi
\[
 \sigma(\Delta)=\operatorname{sgn}(\sigma)\Delta,
\]
gdzie $\operatorname{sgn}:S_n\to\{1,-1\}$
oznacza znak permutacji.

Jeśli $G\subseteq A_n$, wszystkie elementy $G$
mają znak $1$, zatem każdy z nich ustala $\Delta$.
Ponieważ $L/K$ jest skończonym rozszerzeniem Galois,
twierdzenie o ciele elementów stałych daje
$\Delta\in L^G=K$. W takim razie wyróżnik jest
kwadratem elementu $\Delta\in K$.

Odwrotnie, jeśli istnieje $a\in K$ takie, że
$\operatorname{disc}(f)=a^2$, to
$(\Delta-a)(\Delta+a)=0$ w ciele $L$.
Stąd $\Delta=a$ lub $\Delta=-a$, więc $\Delta\in K$.
Dla każdego $\sigma\in G$ mamy jednocześnie
$\sigma(\Delta)=\Delta$ oraz
$\sigma(\Delta)=\operatorname{sgn}(\sigma)\Delta$.
Ponieważ $\Delta\ne0$ i $1\ne-1$ w $K$, koniecznie
$\operatorname{sgn}(\sigma)=1$. Oznacza to
$\sigma\in A_n$ dla każdego $\sigma\in G$.
\end{proof}

\begin{uwaga}[Dlaczego charakterystyka $2$ wymaga osobnego traktowania?]
W charakterystyce $2$ liczby $1$ i $-1$ są równe.
Wobec tego $\Delta$ jest ustalane przez wszystkie
automorfizmy, nawet te wyznaczające nieparzyste
permutacje. Dla rozdzielczego $f$ oznacza to
$\Delta\in K$, a więc $\operatorname{disc}(f)=\Delta^2$
zawsze jest kwadratem w $K$. Nie stanowi to kryterium
zawierania się grupy Galois w $A_n$; poniżej podajemy
konkretny wielomian z grupą $S_3$.
\end{uwaga}

\section{Wielomiany stopnia trzeciego}

\begin{twierdzenie}[Wyróżnik wielomianu sześciennego]\label{aa:tw:cubic-disc}
Niech $K$ będzie dowolnym ciałem, a $p,q\in K$.
Dla $f(X)=X^3+pX+q\in K[X]$ zachodzi
$\operatorname{disc}(f)=-4p^3-27q^2$,
gdzie liczby całkowite $-4$ i $-27$ rozumiemy
jako ich obrazy w $K$ przez homomorfizm $\mathbb Z\to K$.
\end{twierdzenie}
\begin{proof}
Niech $a,b,c\in\overline K$ będą pierwiastkami
$f$, wypisanymi z krotnościami. Ze wzorów Viète'a
mamy $a+b+c=0$, $ab+ac+bc=p$ oraz $abc=-q$.
Kładziemy $u=a+b$ oraz $v=ab$; z pierwszej równości
wynika $c=-u$. Obliczamy kolejno
\[
 p=ab+(a+b)c=v-u^2,\qquad
 q=-abc=uv,
\]
\[
 (a-b)^2=(a+b)^2-4ab=u^2-4v
\]
oraz
\[
 (a-c)(b-c)=(a+u)(b+u)
 =ab+u(a+b)+u^2=v+2u^2.
\]
Zatem, używając definicji wyróżnika,
\[
 \begin{aligned}
 \operatorname{disc}(f)
 &=(a-b)^2(a-c)^2(b-c)^2\\
 &=(u^2-4v)(2u^2+v)^2\\
 &=(u^2-4v)(4u^4+4u^2v+v^2)\\
 &=4u^6-12u^4v-15u^2v^2-4v^3.
 \end{aligned}
\]
Z drugiej strony, podstawiając $p=v-u^2$ i $q=uv$,
dostajemy
\[
 \begin{aligned}
 -4p^3-27q^2
 &=-4(v-u^2)^3-27u^2v^2\\
 &=-4(v^3-3v^2u^2+3vu^4-u^6)-27u^2v^2\\
 &=4u^6-12u^4v-15u^2v^2-4v^3.
 \end{aligned}
\]
Wyrażenia są równe. Każdy krok był tożsamością
wielomianową o współczynnikach całkowitych:
nie dzieliliśmy przez $2$ ani przez $3$.
Dlatego ten sam dowód działa także w charakterystyce
$2$ lub $3$.
\end{proof}

\begin{twierdzenie}[Grupa Galois nierozkładalnego wielomianu sześciennego]\label{aa:tw:cubic-galois}
Niech $K$ będzie ciałem charakterystyki różnej od $2$,
a $f\in K[X]$ monicznym, nierozkładalnym i rozdzielczym
wielomianem stopnia $3$. Wtedy jego grupa Galois jest
izomorficzna z $A_3\cong C_3$, jeśli wyróżnik jest
kwadratem w $K$, a z $S_3$ w przeciwnym razie.
\end{twierdzenie}
\begin{proof}
Niech $L/K$ będzie ciałem rozkładu $f$ i
$G=\operatorname{Gal}(L/K)\hookrightarrow S_3$.
Nierozkładalność $f$ oznacza, że działanie $G$
na jego trzech pierwiastkach jest przechodnie:
każde dwa pierwiastki wielomianu nierozkładalnego
są sprzężone nad $K$, a odpowiednie zanurzenie
$K(\alpha)\to L$ przedłuża się do automorfizmu
normalnego rozszerzenia $L/K$.
Z Twierdzenia~\ref{aa:tw:orbita-stabilizator} wynika, że $3$ dzieli $|G|$.
Z drugiej strony $G$ jest podgrupą $S_3$,
więc na mocy twierdzenia Lagrange'a $|G|$ dzieli
$|S_3|=6$. Jedynymi możliwościami są $|G|=3$ i $6$.
Podgrupa $S_3$ rzędu $3$ jest równa
$A_3=\{\mathrm{id},(123),(132)\}$, a podgrupa rzędu $6$
jest całym $S_3$. Twierdzenie~\ref{aa:tw:disc-alternating}
rozstrzyga między tymi możliwościami na podstawie
wyróżnika, gdy $\operatorname{char}K\ne2$.
\end{proof}

\begin{przyklad}[Dwa możliwe typy]\label{aa:prz:cubic-groups}
Dla $f(X)=X^3-2\in\mathbb Q[X]$ kryterium Eisensteina
przy liczbie pierwszej $2$ daje nierozkładalność;
$f'(X)=3X^2$ nie ma wspólnego pierwiastka z $f$,
więc $f$ jest rozdzielczy. Mamy $p=0$, $q=-2$,
zatem $\operatorname{disc}(f)=-27\cdot4=-108$.
Kwadrat liczby wymiernej nie może być ujemny,
więc grupa Galois jest izomorficzna z $S_3$.

Dla $g(X)=X^3-3X+1$ każdy ewentualny wymierny
pierwiastek musi być całkowitym dzielnikiem $1$,
czyli $1$ albo $-1$. Tymczasem $g(1)=-1$,
$g(-1)=3$, zatem wielomian nie ma pierwiastka
wymiernego i jako sześcienny jest nierozkładalny.
Jego wyróżnik wynosi
$-4(-3)^3-27=108-27=81=9^2\ne0$.
W konsekwencji $g$ jest rozdzielczy i jego grupa
Galois jest cykliczna rzędu $3$.
\end{przyklad}

\begin{przyklad}[Kontrprzykład w charakterystyce dwa]\label{aa:prz:disc-char2}
Zacznijmy od przykładu, w którym grupa jest parzysta.
Nad $\mathbb F_2$ wielomian $X^3+X+1$ nie ma
pierwiastka w $\mathbb F_2$ (wartości w $0$ i $1$
są równe $1$), jest więc nierozkładalny.
Jego ciało rozkładu to $\mathbb F_8$, a grupa Galois
jest cykliczna rzędu $3$, generowana przez Frobeniusa.
Wyróżnik jest równy $-4-27=-31=1$ w $\mathbb F_2$.
Ten przykład sam w sobie \emph{nie} pokazuje błędności
kryterium: grupa rzeczywiście zawiera się w $A_3$.

Teraz niech $K=\mathbb F_2(t)$ będzie ciałem funkcji
wymiernych w zmiennej przestępnej $t$ nad $\mathbb F_2$
i rozważmy $h(X)=X^3+X+t\in K[X]$.
Pokażemy, że $h$ jest rozdzielczy i ma grupę Galois $S_3$,
choć jego wyróżnik jest kwadratem w $K$.

\textbf{Nierozkładalność.}
Gdyby istniał element $r\in K$ spełniający $r^3+r=t$,
to $r\ne0$, ponieważ $t\ne0$. Zapiszmy więc $r=A(t)/B(t)$,
gdzie $A,B\in\mathbb F_2[t]$ są niezerowe
i względnie pierwsze.
Rząd bieguna $r$ w nieskończoności jest równy
$m=\max\{\deg A-\deg B,0\}$.
Jeśli $m=0$, funkcja $r$ nie ma bieguna w nieskończoności,
więc nie ma go też $r^3+r$; tymczasem $t$ ma biegun
rzędu $1$. Jeśli $m>0$, składnik $r^3$ ma biegun rzędu
$3m$, a składnik $r$ rzędu $m$, więc ich suma ma biegun
rzędu $3m\ne1$. Sprzeczność dowodzi, że $h$ nie ma
pierwiastka w $K$, a zatem, jako wielomian stopnia $3$,
jest nierozkładalny.

\textbf{Rozdzielczość.}
Pochodna formalna to $h'(X)=X^2+1=(X+1)^2$.
Wspólny pierwiastek $h$ i $h'$ musiałby być równy $1$,
lecz $h(1)=t\ne0$. Zatem pierwiastki $h$ są różne.

\textbf{Stopień ciała rozkładu.}
Niech $\alpha$ będzie jednym z pierwiastków $h$.
Z równości $t=\alpha^3+\alpha$ wynika
$K(\alpha)=\mathbb F_2(\alpha)$; przy tym $\alpha$
jest przestępny nad $\mathbb F_2$, ponieważ $t$ jest
przestępny. W pierścieniu $K(\alpha)[X]$ mamy
\[
 h(X)=(X+\alpha)
 \bigl(X^2+\alpha X+\alpha^2+1\bigr).
\]
Drugi czynnik ma pierwiastek w $K(\alpha)$
wtedy i tylko wtedy, gdy dla pewnego
$z\in\mathbb F_2(\alpha)$ zachodzi
$z^2+z=1+\alpha^{-2}$: wystarczy w jego równaniu
podstawić $X=\alpha z$ i podzielić przez
$\alpha^2\ne0$ (ponieważ $h(0)=t\ne0$).
Załóżmy, że takie $z$ istnieje. Prawa strona ma
w punkcie $\alpha=0$ biegun rzędu $2$.
Jeśli $z$ nie ma tam bieguna, to $z^2+z$ również
go nie ma. Jeśli $z$ ma biegun rzędu $m>0$,
to $z^2+z$ ma biegun rzędu $2m$. Musi więc być
$m=1$. Zapisując $z$ w postaci ułamka nieskracalnego,
widzimy, że przy biegunie prostym jego część główna
ma postać $c\alpha^{-1}$ dla pewnego
$c\in\mathbb F_2^\times=\{1\}$.
Zatem $z=\alpha^{-1}+w$, gdzie $w$ nie ma bieguna
przy $\alpha=0$. Otrzymujemy
\[
 z^2+z=\alpha^{-2}+\alpha^{-1}+w^2+w,
\]
co nie może być równe $1+\alpha^{-2}$, gdyż
$w^2+w$ nie ma bieguna, a $\alpha^{-1}$ ma.
Sprzeczność dowodzi nierozkładalności czynnika
kwadratowego nad $K(\alpha)$.
Stopień ciała rozkładu $L$ jest zatem
$[L:K]=[L:K(\alpha)][K(\alpha):K]=2\cdot3=6$.
Skoro $L/K$ jest Galois, jego grupa ma rząd $6$
i działa wiernie na trzech pierwiastkach, więc
$\operatorname{Gal}(L/K)\cong S_3$.

Wreszcie Twierdzenie~\ref{aa:tw:cubic-disc},
poprawne także w charakterystyce $2$, daje
$\operatorname{disc}(h)=-4-27t^2=t^2$,
czyli jest on kwadratem niezerowego elementu $t\in K$.
Otrzymaliśmy więc kontrprzykład w charakterystyce $2$:
wyróżnik jest kwadratem w $K$, lecz grupa $S_3$
\emph{nie} zawiera się w $A_3$.
\end{przyklad}
